\documentclass[11pt,a4paper]{article}
\usepackage[utf8]{inputenc}
\usepackage[T1]{fontenc}
\usepackage[spanish]{babel}
\usepackage[margin=2.2cm]{geometry}
\usepackage{amsmath,amssymb}
\usepackage{enumitem}
\usepackage{xcolor}
\usepackage{titlesec}
\usepackage{array,booktabs,longtable}
\usepackage{graphicx}
\usepackage{hyperref}

\hypersetup{colorlinks=true,urlcolor=blue,linkcolor=black}
\setlist{itemsep=2pt,topsep=2pt}
\titlespacing*{\section}{0pt}{10pt}{4pt}
\titlespacing*{\subsection}{0pt}{8pt}{3pt}

\graphicspath{{postproc/}}

\title{\vspace{-1cm}\textbf{TP5 -- Calibración de parámetros numéricos}\\[2pt]\large Kuramoto: elección de $N$, $dt$ y $t_{\text{sim}}$}
\author{Simulación de Sistemas -- 2026Q1}
\date{}

\begin{document}
\maketitle
\vspace{-0.8cm}

Este documento explica \textbf{cómo elegimos} los valores numéricos del TP y \textbf{por qué}, basado en experimentos de convergencia y estacionariedad.

\subsection*{Resultado final}

\begin{center}
\renewcommand{\arraystretch}{1.25}
\begin{tabular}{@{}l l p{8.8cm}@{}}
\toprule
\textbf{Parámetro} & \textbf{Valor} & \textbf{Justificación} \\
\midrule
$N$ & \textbf{600} & Cumple $N>500$ con costo computacional manejable. \\
$dt$ & \textbf{$10^{-3}$} & Test de convergencia RK4: mínimo estable, precisión de máquina. \\
$t_{\text{sim}}$ (completa)  & \textbf{50}   & Sincroniza en $t<1$. Margen $50\times$. \\
$t_{\text{sim}}$ (aleatoria) & \textbf{50}   & Sincroniza en $t<3$ incluso para $p$ chico. \\
$t_{\text{sim}}$ (anillo)    & \textbf{1500} & Transitorios largos por propagación local. \\
\bottomrule
\end{tabular}
\end{center}

\textbf{Costo total estimado del TP:} $\sim$25 min con 4 cores paralelos (vs. $\sim$50 h con $N=1000$, $t_{\text{sim}}=500$ global).

\section{Elección de $N$}

\textbf{Constraint del enunciado:} $N > 500$.

\subsection*{Costo computacional}

Benchmark del motor optimizado (red completa, $K=0.8$, $t_{\text{sim}}=50$):

\begin{center}
\begin{tabular}{c c c}
\toprule
$N$ & Tiempo (s) & Escalado \\
\midrule
500  & 2.92  & $1\times$ \\
750  & 6.67  & $2.3\times$ \\
1000 & 12.32 & $4.2\times$ \\
1500 & 28.37 & $9.7\times$ \\
\bottomrule
\end{tabular}
\end{center}

Costo por paso $\propto N^2$ (cuello de botella: doble bucle de la red completa). El $t_{\text{sim}}$ requerido en anillo escala como $N$ (propagación local), por lo que \textbf{costo total} $\propto N^3$.

\subsection*{Calidad estadística}

Fluctuaciones del parámetro de orden $\sim 1/\sqrt{N}$:
\begin{itemize}
\item $N=600$: $\sim$4.1\%
\item $N=1000$: $\sim$3.2\%
\end{itemize}

Diferencia despreciable. La transición $K_c$ se ve igual de clara en ambos.

\subsection*{Conectividad de la red aleatoria}

Para que un sorteo de Erdős--Rényi sea conectado se necesita $\langle k \rangle > \log(N)$:

\begin{itemize}
\item $N=600$: $\log(600) \approx 6.4 \Rightarrow p_{\min} \approx 0.011$
\item $N=1000$: $\log(1000) \approx 6.9 \Rightarrow p_{\min} \approx 0.007$
\end{itemize}

Con $N=600$ y $p$ chico la red puede estar marginalmente desconectada, mitigado promediando sobre $>10$ realizaciones.

\subsection*{Veredicto: $N = 600$}

Cumple consigna, costo $8\times$ menor que $N=1000$, calidad estadística suficiente.

\section{Elección de $dt$}

\textbf{Setup:} $N=1000$, completa, $K=0.5$, semilla 42, $t_{\text{sim}}=20$. Misma realización corrida con $dt$ variable. Comparamos $r(t)$ tomando $dt=10^{-4}$ como referencia.

\begin{center}
\includegraphics[width=0.96\linewidth]{dt_convergence.png}
\end{center}

\subsection*{Resultados}

\begin{center}
\begin{tabular}{c c l}
\toprule
$dt$ & $\max|\Delta r|$ vs ref & Diagnóstico \\
\midrule
$10^{-1}$    & 0.987            & \textbf{inestable} (catastrófico) \\
$5\cdot10^{-2}$ & 0.981         & inestable \\
$2\cdot10^{-2}$ & 0.977         & inestable \\
$10^{-2}$    & 0.984            & inestable \\
$5\cdot10^{-3}$ & $8.7\cdot10^{-9}$  & converge \\
$2\cdot10^{-3}$ & $3.9\cdot10^{-15}$ & precisión de máquina \\
$\mathbf{10^{-3}}$ & $\mathbf{3.8\cdot10^{-15}}$ & \textbf{precisión de máquina} \\
\bottomrule
\end{tabular}
\end{center}

Hay un \textbf{acantilado de estabilidad} entre $dt = 10^{-2}$ (inestable) y $dt = 5\cdot10^{-3}$ (converge).

\subsection*{Por qué $10^{-3}$ y no $5\cdot10^{-3}$}

\begin{itemize}
\item $5\cdot10^{-3}$ está apenas debajo del threshold; para $K$ mayor o $N$ mayor el threshold puede correrse (condición RK4: $dt \lesssim 2.78/(K\cdot N)$).
\item $10^{-3}$ asegura estabilidad para \textbf{todo el barrido $K\in[0,1]$}.
\item $10^{-3}$ es un valor limpio y conservador.
\end{itemize}

\subsection*{Veredicto: $dt = 10^{-3}$}

\section{Elección de $t_{\text{sim}}$}

Corrimos los casos potencialmente más lentos con $t_{\text{sim}}$ largo (100 $\to$ 500 $\to$ 2000) para ver dónde realmente se estabiliza $r(t)$.

\begin{center}
\includegraphics[width=0.96\linewidth]{tsim_calibration.png}
\end{center}

\subsection*{Por topología}

\textbf{Red completa.} Sincroniza instantáneamente ($t<1$) para cualquier $K>0.05$. $K_c$ efectivo $\approx 2/(\pi\cdot N \cdot g(0)) \approx 10^{-4}$, varios órdenes de magnitud por debajo del rango de interés. $\Rightarrow t_{\text{sim,completa}} = 50$.

\textbf{Red aleatoria.} Incluso con $p=0.01$ sincroniza en $t\approx 2$--3 (diámetro del grafo aleatorio $\sim \log(N)/\log\langle k\rangle$, muy chico). $\Rightarrow t_{\text{sim,aleatoria}} = 50$.

\textbf{Red anillo.} Propagación local: información viaja a sólo $2v$ vecinos, transitorios escalan como $\tau \sim N/v$. Además, para $v$ intermedio ($3$--$7$) y $K$ moderado ($0.5$--$1$) aparecen \textbf{estados quiméricos} y ondas viajeras con transitorios muy largos. Confirmamos comportamiento asintótico a $t_{\text{sim}}=2000$:

\begin{center}
\includegraphics[width=0.96\linewidth]{tsim_n600_long.png}
\end{center}

\begin{center}
\begin{tabular}{l c c c}
\toprule
Caso & $r$ a $t=300$ & $r$ asintótico ($t\approx 1500$) & $\Delta$ \\
\midrule
ring $v=5$, $K=1.0$ & $\sim 0.2$ (mínimo local!) & 0.54 & 0.34 \\
ring $v=5$, $K=0.5$ & $\sim 0.25$ (bajando)      & 0.45 & 0.20 \\
ring $v=10$, $K=0.1$ & $\sim 0.1$ (subiendo)     & 0.37 & 0.27 \\
ring $v=1$, $K=0.5$ & $\sim 0.06$ (oscilante)    & $\sim 0.05$ & OK \\
\bottomrule
\end{tabular}
\end{center}

Si midiéramos $r_\infty$ en $t=300$ daríamos \textbf{valores transitorios incorrectos} para los casos lentos. $\Rightarrow t_{\text{sim,anillo}} = 1500$.

\section{Resumen para writeup}

\begin{quote}
\small
``Calibramos los parámetros numéricos mediante experimentos de convergencia:
\begin{itemize}
\item $dt = 10^{-3}$: elegido tras un test de convergencia con $dt = 10^{-x}$, $x \in \{1,2,3,4\}$ sobre red completa $N=1000$, $K=0.5$. Para $dt \geq 10^{-2}$ el integrador RK4 es inestable; para $dt \leq 10^{-3}$ converge a precisión de máquina.
\item $N = 600$: cumple $N>500$ con costo $\sim 8\times$ menor que $N=1000$ y calidad estadística similar.
\item $t_{\text{sim}}$ por topología: 50 para completa y aleatoria (sincronización en $t<5$), 1500 para anillo (transitorios largos por propagación local, observados en pruebas piloto)."
\end{itemize}
\end{quote}

\end{document}
